\section{Monte Carlo simulation}

This section documents the paper-grade Monte Carlo exercise used to evaluate the contextual policy \textsc{CBI-UCB} with a doubly-robust pseudo-loss under panel missingness and partial/delayed feedback. The accompanying MATLAB package exports a small number of consolidated exhibits (few, but information-dense) into \texttt{out/exhibits/}; the LaTeX below reads those files directly. All reported numbers are Monte Carlo means, with Monte Carlo standard errors (MCSE) in parentheses.
We use the standard UCB exploration bonus (with exploration constant $c$); see \cite{LattimoreSzepesvari2020} for a textbook treatment.

\subsection{Objective}
We study when and why the policy (i) beats the best static arm (value of context) and (ii) fails. The simulation is designed to make the main failure modes referee-visible: sparsity of context\,$\times$\,arm counts due to overly fine discretization, batching and feedback fragmentation, exploration costs, and distortion induced by noisy proxy risk estimates.

\subsection{Data generating process}
For each replication we simulate a panel $Y_{it}$ with a single latent factor,
\begin{equation}
Y_{it} = g(F_t)\,\lambda_i + \varepsilon_{it}, \qquad i=1,\dots,N,\ t=1,\dots,T,
\end{equation}
with $F_t$ following an AR(1) dynamic and $\varepsilon_{it}$ i.i.d. noise. Missingness is drawn cell-by-cell through a logistic rule that depends on an observable contextual covariate built from $(\lambda_i,t)$; this corresponds to a baseline MAR design where selection depends on observables \cite{Rubin1976}, and aligns with propensity-based weighting ideas in missing-data econometrics \cite{Wooldridge2007}. A ``holdout'' mask selects evaluation cells: all MSEs reported below are computed on the holdout.

\paragraph{Econometric motivation.}
The baseline panel model mirrors the dynamic factor structure commonly used in econometric applications with large panels and ragged-edge data releases, where missing values and publication delays are handled via state-space representations and filtering/smoothing methods.\footnote{See e.g. \cite{GiannoneReichlinSmall2008,BanburaModugno2014} for dynamic factor models with staggered releases and arbitrary missingness patterns, and \cite{BaiNg2002} for approximate factor structures in large $N,T$ panels.}

\paragraph{DGP catalog.}
All paper-grade DGPs are declared in the MATLAB file \texttt{src/dgp\_library.m}. The main text focuses on two calibrated environments:
\begin{itemize}
\item \textbf{DFM\_CTX\_MODERATE} (Scenario E; bridge): mild heterogeneity across units, so the value of context is limited and exploration/batching costs can dominate in short horizons.
\item \textbf{DFM\_CTX\_STRONG} (Scenario E2; main): strong heterogeneity (regimes indexed by $\lambda_i$) so the identity of the best imputation arm switches across contexts.
\end{itemize}

Delayed feedback is implemented at the batch level: only holdout cells can become revealed, each with probability $\rho$, and their reveal is delayed by $\Delta$ batches. This creates partial and stale feedback in the pseudo-loss used for learning.

\subsection{Policy, scenarios, and benchmarks}
The action set contains $K=4$ imputation arms. The policy selects one arm per cell using a contextual LCB rule based on doubly-robust risk estimates plus an exploration bonus. We treat:
\begin{itemize}
\item \textbf{Scenario E} as a baseline/bridge environment with simpler/less restrictive feedback (internal Stage 6);
\item \textbf{Scenario E2} as the full setting with delayed/partial feedback (internal Stage 7).
\end{itemize}

Benchmarks are reported alongside the policy:
\begin{itemize}
\item \textbf{Best(static)}: the best arm in hindsight (single arm used everywhere);
\item \textbf{Oracle(ctx)}: the best arm within each context bin (infeasible, but isolates the value of context);
\item \textbf{Policy}: \textsc{CBI-UCB}.
\end{itemize}
We report gaps to both benchmarks, and a sparsity diagnostic for the context\,$\times$\,arm selection counts.

\subsection{Paper-grade design}
The experiment follows the design directives used throughout the paper-grade package:
\begin{enumerate}
\item \textbf{Scaling in $T$ with constant batch length $L$}: we fix $L$ and set $B=\lfloor T/L\rfloor$ so that the feedback delay in \emph{periods} is stable as $T$ grows.
\item \textbf{Context discretization}: the baseline uses a coarse grid ($3\times 3$ bins). An optional 3-level coarsening sensitivity ($3\times 3$, $2\times 2$, $1\times 1$) documents the induced sparsity.
\item \textbf{Exploration sensitivity}: an optional 3-level sensitivity in the UCB constant $c\in\{0.5,1,2\}$ documents exploration costs.
\item \textbf{Scenario bridge (E vs E2)}: the main scaling exhibit reports both scenarios on the same $T$ grid.
\item \textbf{Benchmarks}: Best(static) and Oracle(ctx) are always reported next to the policy.
\item \textbf{MC reporting}: we use $R$ replications and report MCSE for every entry.
\end{enumerate}

\subsection{Exhibits (auto-generated)}
\subsection{Code-generated exhibits}

All tables and figures below are automatically generated by the MATLAB suite and written to
\texttt{out/exhibits/}. To refresh them, run \texttt{RUN\_ALL.m} (one click) or the driver
\texttt{S20\_main\_T\_scaling\_stage7.m} (main exhibits only).

\paragraph{Main scaling exhibit.}
\begin{table}[!ht]
\centering
\caption{Scaling in $T$ with constant batch length $L$. Entries are Monte Carlo means with MCSE in parentheses.
We report Best(static), Policy, Oracle(ctx), the excess risk $\text{Gap}=\text{Policy}-\text{Best(static)}$,
context value $\text{CtxGain}=\text{Best(static)}-\text{Oracle(ctx)}$, oracle regret
$\text{Regret}=\text{Policy}-\text{Oracle(ctx)}$, win rate $\mathbb{P}(\text{Gap}<0)$, and a sparsity diagnostic
for context$\times$arm counts. Since $\text{Gap}=\text{Regret}-\text{CtxGain}$, the win rate equals
$\mathbb{P}(\text{Regret}<\text{CtxGain})$.}
\vspace{0.25em}
\IfFileExists{\ExhibitDir tab_paper_scaling.tex}{
  \input{\ExhibitDir tab_paper_scaling.tex}
}{
  \fbox{\parbox{0.9\linewidth}{Missing file: \texttt{tab\_paper\_scaling.tex}. Run \texttt{S20\_main\_T\_scaling\_stage7.m}.}}
}
\end{table}

\begin{figure}[!ht]
\centering
\caption{Performance as a function of $T$ under Scenario E (baseline) and Scenario E2 (partial/delayed feedback).
Left panels plot oracle regret $\text{Regret}=\text{Policy}-\text{Oracle(ctx)}$; right panels plot win rate
$\mathbb{P}(\text{Gap}<0)$. Using $\text{Gap}=\text{Regret}-\text{CtxGain}$ with
$\text{CtxGain}=\text{Best(static)}-\text{Oracle(ctx)}$, the win rate equals $\mathbb{P}(\text{Regret}<\text{CtxGain})$.
Error bars are $\pm 2$ MCSE.}
\vspace{0.25em}
\IfFileExists{\ExhibitDir fig_paper_scaling.pdf}{
  \includegraphics[width=0.98\linewidth]{fig_paper_scaling.pdf}
}{
  \fbox{\parbox{0.9\linewidth}{Missing file: \texttt{fig\_paper\_scaling.pdf}. Run \texttt{S20\_main\_T\_scaling\_stage7.m}.}}
}
\end{figure}

\ifShowCalibration
\paragraph{Sensitivity (optional).}
\begin{table}[!ht]
\centering
\caption{Sensitivity around $(\#\text{bins},c)$ at the calibration $T$ (Monte Carlo means with MCSE in parentheses).}
\vspace{0.25em}
\IfFileExists{\ExhibitDir tab_paper_calibration.tex}{
  \input{\ExhibitDir tab_paper_calibration.tex}
}{
  \fbox{\parbox{0.9\linewidth}{Missing file: \texttt{tab\_paper\_calibration.tex}.}}
}
\end{table}
\fi


\subsection{How to read the main exhibit (interpretation guide)}
The main exhibit is organized around the same object a referee expects in simulation work: a \emph{functional} of performance as the horizon grows.
The figure reports two complementary views of the policy performance as $T$ increases (holding $L$ fixed):
\begin{itemize}
\item \emph{oracle regret} (Policy MSE minus Oracle(ctx) MSE), which should decrease toward zero as the policy learns the best arm in each context;
\item \emph{win rate} $\mathbb P(\text{Gap}<0)$, i.e. the fraction of replications in which the policy beats Best(static) (a power-like diagnostic).
\end{itemize}
\noindent
To interpret these objects jointly, decompose the excess risk against the best static arm as
\begin{equation}
\text{Gap} \,=\, \text{Policy}-\text{Best(static)} \,=\, \underbrace{\bigl(\text{Policy}-\text{Oracle(ctx)}\bigr)}_{\text{Regret}} \, - \,
\underbrace{\bigl(\text{Best(static)}-\text{Oracle(ctx)}\bigr)}_{\text{CtxGain}}.
\end{equation}
Hence the win rate can be read as a ``threshold crossing'' probability,
\begin{equation}
\mathbb P(\text{Gap}<0) \,=\, \mathbb P\bigl(\text{Regret}<\text{CtxGain}\bigr),
\end{equation}
so failures can be attributed either to low context value (small CtxGain) or to slow/biased learning (large regret).
Sparsity of context$\times$arm counts is reported in the \emph{table} to make failures attributable to discretization/feedback fragmentation rather than to ``bad luck''. (Optional sensitivity tables can be enabled if we want to further isolate the bins--sparsity trade-off and exploration costs.)

\subsection{How to reproduce}
From MATLAB:
\begin{enumerate}
\item optionally run \texttt{WIPE\_OUTPUTS.m} (clears \texttt{out/});
\item run \texttt{S20\_main\_T\_scaling\_stage7.m} (writes the main exhibits and \texttt{run\_manifest.txt});
\item optionally run \texttt{S10\_calibrate\_bins.m} and \texttt{S11\_calibrate\_ucb\_c.m} (writes a sensitivity table in \texttt{out/exhibits/});
\item compile \texttt{text/PaperMCMain.tex} (or use the one-click driver \texttt{RUN\_ALL.m}).
\end{enumerate}
